% Generated draft: requires mathematical and source review.
\documentclass[11pt]{article}
\usepackage[margin=25mm]{geometry}
\usepackage{fontspec}
\setmainfont{Latin Modern Roman}
\usepackage{amsmath,amssymb,mathtools,unicode-math}
\setmathfont{Latin Modern Math}
\usepackage{enumitem,needspace,xurl,hyperref}
\setlength{\parindent}{0pt}
\setlength{\parskip}{5pt}
\setlength{\emergencystretch}{3em}
\newcommand{\N}{\mathbb{N}}
\newcommand{\Z}{\mathbb{Z}}
\newcommand{\Q}{\mathbb{Q}}
\newcommand{\R}{\mathbb{R}}
\newcommand{\C}{\mathbb{C}}
\newcommand{\F}{\mathbb{F}}
\newcommand{\A}{\mathbb{A}}
\newcommand{\PP}{\mathbb{P}}
\newcommand{\E}{\mathbb{E}}
\newcommand{\eps}{\varepsilon}
\newcommand{\dd}{\,\mathrm{d}}
\begin{document}
\section*{UnsolvedMath ID: 30005542}
\textbf{Source problem code:} catalogue entry
\subsection{Definitions and mathematical statement}
\textbf{Unrecoverable mathematical statement.} The supplied record for permanent UnsolvedMath ID \(30005542\) contains a table-of-contents fragment, not a mathematical problem. Its opening word ``problem'' is followed by a page reference to 1405 without an identifiable associated statement. The remaining entries give talk titles, contributors, and page references from 1407 through 1425. These entries do not supply definitions, hypotheses, domains, boundary conditions, quantifiers, or a conclusion to be proved or disproved. Consequently, no precise mathematical statement or associated definitions can be recovered from the supplied evidence.

The registry title, ``Unspecified Harmonic Analysis Incidence Problem,'' does not identify an intended mathematical target. None of the listed talks can be selected as that target from this record. The supplied \texttt{statement\_status} is \texttt{unrecoverable}, and the supplied verification states that no separate clean statement has been promoted; the original record remains publicly visible for provenance. The registry's \texttt{open} label reflects the supplied curation's classification of an unclear resolution, rather than a verified open conjecture.
\subsection{Short English statement}
This record preserves a workshop table-of-contents fragment. It contains no recoverable mathematical question. A substantive source passage identifying the intended problem is needed before definitions and a precise statement can be written. The supplied ``open'' classification is a curation label, not an independently verified mathematical status.
\subsection{Sources}
\begin{itemize}
\item Local UnsolvedMath registry, permanent ID \(30005542\), exact source problem code \texttt{catalogue entry}, in the collection \emph{Oberwolfach Reports Open Problems}. The record supplies proposed year 2023 but no proposer and no clean statement. Its original and displayed statements contain the same table-of-contents fragment, retained as context and provenance.
\item Supplied report citation: \emph{Incidence Problems in Harmonic Analysis, Geometric Measure Theory, and Ergodic Theory} (2023), Oberwolfach Reports, pp.~1397--1452, DOI \texttt{10.4171/owr/2023/25}. Supplied link: \url{https://doi.org/10.4171/owr/2023/25}. The supplied background describes Oberwolfach Reports as documenting mathematical workshops and their open-problem discussions; this does not identify a problem in the extracted fragment.
\item Supplied literature source, labeled ``OWR 2023 25'': \url{https://publications.mfo.de/bitstream/handle/mfo/4055/OWR_2023_25.pdf?sequence=1}. The literature assessment dated \textbf{2026-08-22} is supplied dated curation, not independently verified here. It classifies the record as open because the extracted value is a table-of-contents fragment rather than a well-defined mathematical claim, leaving its resolution unclear. The supplied background explicitly calls this dated literature triage and requests independent verification of time-sensitive status and cited claims.
\item Attribution within the supplied fragment, preserved as contents entries rather than attribution of a conjecture:
\begin{itemize}
\item Unidentified preceding entry ending with ``problem,'' p.~1405.
\item Balázs Bárány, joint with Roope Anttila and Antti Käenmäki: ``On the Assouad dimension of the Takagi function,'' p.~1407.
\item Amir Algom, joint with Yuanyang Chang, Meng Wu and Yu-Liang Wu: ``Polynomial Fourier decay for smooth images of self-similar measures,'' p.~1410.
\item Marianna Csörnyei: ``Effective methods in geometric measure theory,'' p.~1411.
\item Alexia Yavicoli: ``Thickness and a Gap Lemma in Rd,'' p.~1412. The source's notation ``Rd'' is retained without interpreting it as a specified mathematical domain.
\item Shaoming Guo: ``Orthogonal projections of fractal sets,'' p.~1413.
\item Lars Becker: ``Sharp Fourier extension from the circle for functions with localized support,'' p.~1414.
\item Contributor rendered in the extraction as ``Damian Da˛ browski'': ``Quantifying Besicovitch projection theorem,'' p.~1415.
\item Marina Iliopoulou, joint with Anthony Carbery and Bassam Shayya: ``Some remarks on the Mizohata--Takeuchi conjecture,'' p.~1417.
\item Tamás Keleti, joint with András Máthé: ``Equivalences between different forms of the Kakeya conjecture and duality of Hausdorff and packing dimensions for additive complements,'' p.~1420.
\item Orit E. Raz, joint with Joshua Zahl: ``Hausdorff-dimension analogue of the Elekes--Ronyai theorem and applications,'' p.~1422; the supplied extraction renders the theorem name as ``Elekes--Ronyái.''
\item Michael Hochman: ``Symmetries of Cantor sets,'' p.~1423.
\item Lauritz Streck: ``Entropy Rates in Fractal Geometry,'' p.~1425.
\end{itemize}
\end{itemize}
\end{document}
